\documentclass[border=4pt]{standalone}
\usepackage{amsmath,amssymb,mathtools,xcolor,varwidth}
\definecolor{corrorange}{HTML}{E8770A}
\definecolor{figaviolet}{HTML}{8E24AA}
\definecolor{figbteal}{HTML}{00897B}
\begin{document}
\begin{varwidth}{28em}
\large\bfseries
The \textcolor{corrorange}{correspondence lines} pair each strip of the upper figure with its mate below: $p_1$ with $q_1$, $p_2$ with $q_2$, $p_3$ with $q_3$. By hypothesis every such pair shares one ultimate ratio $k$, so that $p_i / q_i \to k$ each to each respectively. Then, by composition, the sum of \textcolor{figaviolet}{all the upper parallelograms} is to the sum of \textcolor{figbteal}{all the lower parallelograms} in that same ratio $k$, since adding terms of equal ratio preserves the ratio. Finally, by Lemma III each rank of parallelograms is ultimately equal to its own curvilinear figure --- so the sums may be replaced by the whole figures, giving $AacE : PprT = k$. Thus two figures whose corresponding strips share an ultimate ratio are themselves in that very ratio. \textit{Q.E.D.}
\end{varwidth}
\end{document}
